\chapter{1961年全国卷}

\section{解答题}

\begin{problem}
\begin{enumerate}
\item 求 \(\left(2-x\right)^{10}\) 展开式里 \(x^7\) 的系数.

\item 解方程 \(2\lg x=\lg\left(x+12\right)\). （注：\(\lg\) 是以 \(10\) 为底的对数的符号.）

\item 求函数 \(y=\frac{\sqrt{x-1}}{x-5}\) 的自变量 \(x\) 的允许值的范围.

\item 求 \(\sin\frac{\pi}{12}\sin\frac{5\pi}{12}\) 的值.

\item 一个水平放着的圆柱形排水管，内半径是 \(12\text{ 厘米}\). 排水管的圆截面上被水淹没部分的弧含有 \(150^\circ\)（如图）. 求这个截面上有水部分的面积. （取 \(\pi=3.14\)）

\item 已知 \(\triangle ABC\) 的一边 \(BC\) 在平面 \(N\) 内，\(\triangle ABC\) 所在的平面与平面 \(N\) 组成二面角 \(a\)（\(a<90^\circ\)）. 从 \(A\) 点作平面 \(N\) 的垂线 \(AA'\)，\(A'\) 是垂足. 设 \(\triangle ABC\) 的面积是 \(S\)，求证 \(\triangle A'BC\) 的面积是 \(S\cos a\).
\end{enumerate}

\FigureLayoutDeclare{img/1961/national/q1_5_fig1.png}{5.5cm}{0bp 0bp 0bp 0bp}
\bitmapfigure[max width=6.0cm]{img/1961/national/q1_5_fig1.png}
\end{problem}

\begin{solution}
\begin{enumerate}
\item 二项式展开式的通项为 \(\mathrm C_{10}^{r}2^{10-r}(-x)^r\)，其中 \(r=0\)，\(1\)，\(\ldots\)，\(10\). 要得到 \(x^7\)，必须取 \(r=7\)，所以所求系数为
\[
\mathrm C_{10}^{7}2^3(-1)^7=-\mathrm C_{10}^{3}\cdot8=-\frac{10\cdot9\cdot8}{1\cdot2\cdot3}\cdot8=-960
\]

\item 由对数的定义，原方程有意义必须满足 \(x>0\)，且此时 \(x+12>0\) 自动成立. 在此条件下，\(2\lg x=\lg x^2\)，故原方程等价于 \(\lg x^2=\lg\left(x+12\right)\)，从而 \(x^2=x+12\). 因此
\(x^2-x-12=0\)，\(\left(x-4\right)\left(x+3\right)=0\).
解得 \(x=4\) 或 \(x=-3\). 其中 \(x=-3\) 不满足 \(x>0\)，代回原方程也使 \(\lg x\) 无意义，舍去；\(x=4\) 满足定义域且代入原方程成立. 所以原方程的解为 \(x=4\).

\item 根式有意义要求 \(x-1\geq0\)，即 \(x\geq1\). 分母不能为零，要求 \(x-5\neq0\)，即 \(x\neq5\). 因此自变量 \(x\) 的允许值的范围为 \(x\in\left[1,5\right)\cup\left(5,+\infty\right)\).

\item 因为 \(\frac{5\pi}{12}=\frac{\pi}{2}-\frac{\pi}{12}\)，所以 \(\sin\frac{5\pi}{12}=\cos\frac{\pi}{12}\). 从而
\[
\sin\frac{\pi}{12}\sin\frac{5\pi}{12}=\sin\frac{\pi}{12}\cos\frac{\pi}{12}=\frac{1}{2}\sin\frac{\pi}{6}=\frac{1}{4}
\]

也可以利用积化和差公式：
\[
\sin\frac{\pi}{12}\sin\frac{5\pi}{12}=\frac{1}{2}\left(\cos\frac{\pi}{3}-\cos\frac{\pi}{2}\right)=\frac{1}{2}\left(\frac{1}{2}-0\right)=\frac{1}{4}
\]

\item 设圆心为 \(O\)，水面弦的端点为 \(A\)、\(B\)，水淹没的弧为弧 \(ACB\). 因为弧 \(ACB\) 所对的圆心角为 \(150^\circ\)，所以扇形 \(O\text{-}ACB\) 的面积为
\[
\frac{150}{360}\pi\cdot12^2=60\pi\text{ 平方厘米}
\]
三角形 \(OAB\) 的两边 \(OA=OB=12\)，夹角为 \(150^\circ\)，所以
\[
S_{\triangle OAB}=\frac{1}{2}\cdot12^2\sin150^\circ=36\text{ 平方厘米}
\]
水面以下的部分是扇形 \(O\text{-}ACB\) 减去三角形 \(OAB\)，故其面积为
\[
60\pi-36=60\times3.14-36=152.4\text{ 平方厘米}
\]
所以这个截面上有水部分的面积是 \(152.4\text{ 平方厘米}\).

\item 作 \(AD\perp BC\) 于 \(D\)，并连接 \(A'D\). 由于 \(AD\) 在平面 \(N\) 上的正投影为 \(A'D\)，而 \(AD\perp BC\)，且 \(BC\) 是平面 \(N\) 与平面 \(ABC\) 的交线，所以 \(A'D\perp BC\). 因而由两平面分别作出的垂直于交线的直线可知，\(\angle ADA'\) 是二面角的平面角，故 \(\angle ADA'=a\).

\solutionfigure{\bitmapfigure[max width=4.2cm]{img/1961/national/q1_6_solution_fig1.png}}

因为 \(AA'\perp\) 平面 \(N\)，所以 \(AA'\perp A'D\)，三角形 \(AA'D\) 是直角三角形，且 \(AD\) 是斜边. 因此 \(A'D=AD\cos a\). 又 \(AD\) 是三角形 \(ABC\) 对底边 \(BC\) 的高，所以
\[
S_{\triangle A'BC}=\frac{1}{2}BC\cdot A'D=\frac{1}{2}BC\cdot AD\cos a=S\cos a
\]
故三角形 \(A'BC\) 的面积为 \(S\cos a\).
\end{enumerate}
\end{solution}

\begin{problem}
一个机器制造厂的三年生产计划，每年比上一年增产机器的台数相同. 如果第三年比原计划多生产 \(1\text{ 千台}\)，那末每年比上一年增长的百分数就相同，而且第三年生产的台数恰等于原计划三年生产总台数的一半，原计划每年各生产机器多少台？
\end{problem}

\begin{solution}
解法一：设原计划第一、二、三年生产机器的台数分别为 \(x\text{ 千台}\)、\((x+y)\text{ 千台}\)、\((x+2y)\text{ 千台}\)，则新计划第三年为 \((x+2y+1)\text{ 千台}\). 原计划每年增加 \(y\text{ 千台}\)，因此新计划第一年到第二年的增长率为 \(\frac{y}{x}\)，第二年到第三年的增长率为 \(\frac{y+1}{x+y}\). 题意给出这两个增长率相同，所以
\[
\frac{y}{x}=\frac{y+1}{x+y}
\]
交叉相乘得 \(y(x+y)=x(y+1)\)，即 \(y^2=x\). 又因为新计划第三年产量等于原计划三年总产量的一半，所以
\[
x+2y+1=\frac{1}{2}\left[x+(x+y)+(x+2y)\right]
\]
化简得 \(x-y-2=0\). 因而需要解
\[
\begin{cases}
y^2-x=0,\\
x-y-2=0.
\end{cases}
\]
由 \(x=y+2\)，代入第一式得 \(y^2-y-2=0\)，即 \((y-2)(y+1)=0\). 所以 \(y=2\) 或 \(y=-1\). 当 \(y=2\) 时，\(x=4\)，三年的产量依次为 \(4\text{ 千台}\)、\(6\text{ 千台}\)、\(8\text{ 千台}\)；当 \(y=-1\) 时，三年的产量为 \(1\text{ 千台}\)、\(0\text{ 千台}\)、\(-1\text{ 千台}\)，且增长率中的分母 \(x+y=0\)，不符合生产计划的实际意义，舍去.

解法二：设原计划第一、二、三年生产机器的台数分别为 \((x-y)\text{ 千台}\)、\(x\text{ 千台}\)、\((x+y)\text{ 千台}\)，则新计划第三年为 \((x+y+1)\text{ 千台}\). 由相邻两年的增长率相同，得
\[
\frac{y}{x-y}=\frac{y+1}{x}
\]
即 \(xy=(x-y)(y+1)\)，化简为 \(y^2+y-x=0\). 又由总产量条件，得
\[
x+y+1=\frac{1}{2}\left[(x-y)+x+(x+y)\right]
\]
所以 \(x-2y-2=0\). 于是
\[
\begin{cases}
y^2+y-x=0,\\
x-2y-2=0.
\end{cases}
\]
由 \(x=2y+2\)，代入第一式得 \(y^2-y-2=0\)，所以 \(y=2\) 或 \(y=-1\). 当 \(y=2\) 时，\(x=6\)，原计划产量为 \(4\text{ 千台}\)、\(6\text{ 千台}\)、\(8\text{ 千台}\)；当 \(y=-1\) 时，原计划第二、三年产量不符合题意，舍去. 因此原计划第一、二、三年各生产机器 \(4\text{ 千台}\)、\(6\text{ 千台}\)、\(8\text{ 千台}\).
\end{solution}

\begin{problem}
有一块圆环形的铁皮，它的内半径是 \(45\text{ 厘米}\)，外半径是 \(75\text{ 厘米}\). 用它的五分之一（如图中的阴影部分），作圆台形水桶的侧面，这个水桶的容积是多少立方厘米？

\FigureLayoutDeclare{img/1961/national/q3_fig1.png}{5.0cm}{0bp 0bp 0bp 0bp}
\bitmapfigure[max width=6.0cm]{img/1961/national/q3_fig1.png}
\end{problem}

\begin{solution}
将圆环形铁皮的五分之一卷成圆台侧面时，内、外圆弧分别成为水桶下底和上底的圆周. 因此水桶下底半径为
\[
r=\frac{1}{5}\times45=9\text{ 厘米}
\]
上底半径为
\[
R=\frac{1}{5}\times75=15\text{ 厘米}
\]
圆台的母线就是圆环的宽度，所以母线长为 \(75-45=30\text{ 厘米}\). 过圆台轴作截面，得到直角三角形，其斜边为母线，另一直角边为两底半径之差 \(R-r=6\text{ 厘米}\). 因此圆台的高为
\[
h=\sqrt{30^2-6^2}=\sqrt{864}=12\sqrt6\text{ 厘米}
\]
圆台容积为
\[
V=\frac{1}{3}\pi h\left(R^2+Rr+r^2\right)=\frac{1}{3}\pi\cdot12\sqrt6\left(15^2+15\times9+9^2\right)=1764\sqrt6\pi\text{ 立方厘米}
\]
所以水桶的容积是 \(1764\sqrt6\pi\text{ 立方厘米}\).
\end{solution}

\begin{problem}
在平地上有 \(A\)、\(B\) 两点. \(A\) 在山的正东. \(B\) 在山的东南，且在 \(A\) 的西 \(65^\circ\) 南 \(300\text{ 米}\) 的地方. 在 \(A\) 测得山顶的仰角是 \(30^\circ\)，求山高. （\(\sin70^\circ=0.940\)，精确到 \(10\text{ 米}\).）
\end{problem}

\begin{solution}
设山顶为 \(D\)，山脚在平地上的垂足为 \(C\)，则 \(CD\) 是山高. 由 \(A\) 在山的正东、\(B\) 在山的东南，得地面三角形 \(ABC\) 中 \(\angle ACB=45^\circ\). 又由 \(B\) 在 \(A\) 的西 \(65^\circ\) 南，得 \(\angle CAB=65^\circ\). 因而
\[
\angle ABC=180^\circ-45^\circ-65^\circ=70^\circ
\]
在三角形 \(ABC\) 中，已知 \(AB=300\text{ 米}\)，由正弦定理
\[
\frac{AC}{\sin70^\circ}=\frac{AB}{\sin45^\circ}
\]
所以 \(AC=\frac{300\sin70^\circ}{\sin45^\circ}=300\sqrt2\sin70^\circ\text{ 米}\).
在直角三角形 \(ACD\) 中，\(\angle CAD=30^\circ\)，所以
\[
CD=AC\tan30^\circ=300\sqrt2\sin70^\circ\cdot\frac{1}{\sqrt3}=100\sqrt6\sin70^\circ\text{ 米}
\]
代入 \(\sin70^\circ=0.940\)，并取 \(\sqrt6\approx2.45\)，得 \(CD\approx100\times2.45\times0.940=230.3\text{ 米}\). 按题意精确到 \(10\text{ 米}\)，山高为 \(230\text{ 米}\).

\solutionfigure{\bitmapfigure[max width=4.2cm]{img/1961/national/q4_solution_fig1.png}}
\end{solution}

\section{选作题}

\begin{problem}
\(k\) 是什么实数的时候，方程 \(x^2-2\left(k+3\right)x+3k^2+1=0\) 有实数根？
\end{problem}

\begin{solution}
因为二次项系数为 \(1\)，方程有实数根的充要条件是判别式不小于零. 这里
\[
\Delta=\left[-2\left(k+3\right)\right]^2-4\left(3k^2+1\right)
=4\left[\left(k+3\right)^2-\left(3k^2+1\right)\right]
=4\left(-2k^2+6k+8\right)
\]
所以 \(\Delta\geq0\) 等价于 \(-2k^2+6k+8\geq0\). 两边同除以 \(-2\) 时不等号方向改变，得
\(k^2-3k-4\leq0\)，\(\left(k+1\right)\left(k-4\right)\leq0\).
因此 \(-1\leq k\leq4\). 当 \(k=-1\) 或 \(k=4\) 时判别式为零，有重实根；当 \(-1<k<4\) 时判别式为正，有两个不相等的实根. 所以所求实数范围为 \(\left[-1,4\right]\).
\end{solution}

\begin{problem}
设方程 \(8x^2-\left(8\sin a\right)x+2+\cos2a=0\) 的两个根相等，求 \(a\).
\end{problem}

\begin{solution}
该方程是关于 \(x\) 的二次方程，且二次项系数 \(8\neq0\). 两根相等的充要条件是判别式为零，因此
\[
\left(8\sin a\right)^2-4\cdot8\left(2+\cos2a\right)=0
\]
两边除以 \(32\)，得 \(2\sin^2a-\left(2+\cos2a\right)=0\). 利用 \(\cos2a=1-2\sin^2a\)，可得
\(2\sin^2a-\left(3-2\sin^2a\right)=0\)，\(4\sin^2a-3=0\).
所以 \(\sin a=\pm\frac{\sqrt3}{2}\).

当 \(\sin a=\frac{\sqrt3}{2}\) 时，\(a=n\pi+\left(-1\right)^n\frac{\pi}{3}\)，其中 \(n\in\Z\)；当 \(\sin a=-\frac{\sqrt3}{2}\) 时，\(a=n\pi-\left(-1\right)^n\frac{\pi}{3}\)，其中 \(n\in\Z\). 合并两类解，得到
\(a=n\pi\pm\frac{\pi}{3}\)，\(n\in\Z\).
\end{solution}
